\chapter{Discusión}
\label{cap:discusion}

\section{Ventajas y límites de cada método}

\begin{table}[ht]
  \centering
  \caption{Ventajas y limitaciones observadas en la reproducción local.}
  \label{tab:ventajas}
  \small
  \begin{tabular}{p{0.30\textwidth}p{0.32\textwidth}p{0.30\textwidth}}
    \toprule
    Método & Ventaja principal & Límite principal \\
    \midrule
    Baseline XGBoost CPU &
      Mínimo costo de desarrollo; piso de referencia inmediato ($\sim$0.92). &
      Features crudas desperdician la señal de entidad; 112\,s por ajuste frena la iteración. \\
    \addlinespace
    Baseline XGBoost GPU &
      4.4$\times$ más rápido con AUC idéntico; habilita búsqueda de hiperparámetros. &
      No aporta precisión por sí solo; exige VRAM suficiente. \\
    \addlinespace
    LightGBM + RFE + count encoding (nroman) &
      30\,\% menos features con AUC equivalente; validación temporal estricta; interpretable. &
      Ganancia de precisión marginal; la tasa de aprendizaje 0.0069 exige miles de árboles. \\
    \addlinespace
    XGBoost d12 + magic UID (cdeotte) &
      Mayor AUC local (0.9479); ganancia atribuible y replicable ($+0.0104$). &
      Requiere conocimiento del dominio (composición del UID); agregaciones costosas; protocolo
      transductivo (ver \cref{sec:limites}). \\
    \bottomrule
  \end{tabular}
\end{table}

\section{Qué importa realmente: jerarquía de impacto}

Los experimentos permiten ordenar los factores por su contribución al AUC:

\begin{enumerate}
  \item \textbf{Validación honesta} (no aporta AUC, evita ilusiones): el esquema de validación
        decide si un resultado es real. KFold aleatorio habría reportado valores notablemente
        superiores e irrepetibles en producción.
  \item \textbf{Ingeniería de entidades} ($+0{,}028$ acumulado): el pseudo-cliente y sus
        agregaciones explican la mayor parte de la distancia entre el baseline y el tope local.
        La lección es transferible a cualquier dominio con entidades implícitas (usuarios,
        dispositivos, comercios): \emph{construir la entidad, agregar su historial, descartar la
        llave}.
  \item \textbf{Estabilización temporal} (habilitadora): normalizar columnas \texttt{D} y filtrar
        por consistencia temporal no sube el AUC por sí sola, pero hace que las ganancias
        anteriores \emph{se transfieran} al futuro.
  \item \textbf{Selección de variables} ($\approx$neutral aquí): recortar de 432 a 120 variables
        mantuvo el AUC con menos costo, relevante en producción.
  \item \textbf{Algoritmo e hiperparámetros} (menor): entre XGBoost y LightGBM la diferencia
        observada es inferior al efecto de cualquier técnica de datos; la optimización bayesiana
        de nroman vale menos que una sola feature de entidad bien construida.
  \item \textbf{Hardware} ($0$ AUC, $\times$4.4 velocidad): la GPU acelera el ciclo de
        experimentación, que es donde se generan las ganancias de datos.
\end{enumerate}

\section{Limitaciones metodológicas}
\label{sec:limites}

\begin{enumerate}
  \item \textbf{Protocolo transductivo}: las agregaciones por UID se computan sobre todo el
        período de entrenamiento, incluidos los meses usados como validación (sin usar sus
        etiquetas). El AUC local absoluto es, por tanto, un límite superior respecto de un
        sistema causal estricto. La variante estricta (estadísticas solo con transacciones
        anteriores a cada transacción) es el siguiente hito propuesto.
  \item \textbf{Un solo grano de validación para los XGB d12}: GroupKFold de 3 meses (6 en el
        original) por restricciones de tiempo de GPU; las diferencias entre folds fueron
        pequeñas y consistentes en dirección.
  \item \textbf{Corpus no ejecutado}: 8 notebooks de alto valor analítico (EDA, re-muestreo,
        optimización bayesiana completa, \emph{blend} de artgor, RAPIDS de Deotte) se analizaron
        cualitativamente pero no se reprodujeron numéricamente; sus conclusiones no cuentan con
        medición local.
  \item \textbf{Métrica única}: el AUC ordena bien pero no responde por costos asimétricos de
        negocio (falso positivo bloquea un pago legítimo); un despliegue real requiere umbral y
        métrica de utilidad propios.
\end{enumerate}

\section{Amenazas a la validez}

Las adaptaciones de librerías (XGBoost $\ge$2, LightGBM $\ge$4, pandas~3) alteran detalles de
desempate y de parada temprana respecto de 2019; se priorizó preservar la \emph{lógica} de cada
pipeline sobre la igualdad bit a bit. Las diferencias con los puntajes históricos de Kaggle son
esperadas y se documentan en \cref{cap:resultados}; ninguna cambia la conclusión comparativa.
